\begin{solapartat}
\begin{minipage}[c]{0.5\linewidth}
\punts{0,5} $\left\{\begin{array}{l}
E = \dfrac{\Delta V}{d}\\[2ex]
E = \dfrac{12}{0,06} = 200\ \mathrm{V/m}
\end{array}\right.$

\medskip
\punts{0,2} Camp elèctric dirigit cap a l'esquerra.
\end{minipage}\hfill
\begin{minipage}[c]{0.45\linewidth}
\figura[0.85]{sol_fig1.png}
\centering\punts{0,3}
\end{minipage}

\medskip
Els estudiants poden indicar un nombre diferent de superfícies equipotencials, superior a 1, respectant l'equidistància.
\end{solapartat}

\begin{solapartat}
\punts{0,2} $W_{ext} = q\Delta V$

\punts{0,6} $\Delta V = 12 - 7 = 5\ \mathrm{V}$

\punts{0,2} $W_{ext} = 0,1\cdot 10^{-6}\cdot 5 = 5\cdot 10^{-7}\ \mathrm{J}$
\end{solapartat}
